We have demonstrated a robust and reproducible method for generating millimeter-scale antibubbles by replacing the random foam layer of traditional foam-based techniques with a 3D-printed triangular-prism frame that fixes a single Plateau border in space. This geometric simplification enables systematic variation of droplet radius (R = 1.0-2.2 mm), impact velocity (v0 = 0.1-1.8 m/s), and ambient pressure (P = 0.1-1.0 atm) across 199 controlled experiments, revealing the underlying physics of antibubble formation.

High-speed imaging identifies four universal stages of droplet-PB interaction -- free fall, entering the PB, transport, and release -- and three outcomes: packing (successful multi-layer formation), outer coalescence, and inner coalescence. From these observations, we derive two dimensionless criteria that together predict the parameter space for antibubble generation. The We-Bo criterion [Eq. (1), We + 2 Bo = 0.038 +/- 0.004], obtained from energy conservation, defines the minimum energy required to deform the PB films and form the encapsulating 1G1L structure. The t/t criterion [Eq. (3), t_p/t_d < 1], derived from lubrication theory applied to the gas film between the droplet and the deforming soap films, ensures that pinch-off completes before the gas film drains. Together, 97% of packing outcomes fall within the region satisfying both criteria.

Two findings have broader implications. First, the PB method's energy criterion constant is roughly half that of horizontal-interface methods, providing a quantitative explanation for why the inclined PB geometry achieves ~80% yield while horizontal methods remain at ~20%. The concave, pre-entrained air pocket at the PB node fundamentally lowers the energy barrier for gas-film encapsulation. Second, the t/t criterion reveals a pressure-dependent failure mode: at reduced ambient pressure (P < 0.3 atm), the decreased gas density prolongs gas-film drainage, causing inner coalescence even when the We-Bo energy condition is met. This result carries practical consequences for space and high-altitude applications where reduced ambient pressure may destabilize antibubble formation, and it suggests that the operating pressure must be explicitly accounted for in any antibubble-generation protocol.

The fixed-PB platform demonstrated here opens several avenues. The deterministic PB geometry enables controlled studies of droplet transport, coalescence, and gas-film stability in a well-characterized microchannel -- a simplified physical model of the complex foam networks encountered in industrial and biological settings. The platform can be extended to rigid-particle transport (as our control experiments with PP, PMMA, PS, and POM spheres demonstrate) to disentangle droplet-specific phenomena such as shape oscillation from universal confined-transport dynamics. Finally, the vacuum-chamber capability directly prepares for planned microgravity experiments, where antibubble lifetimes become virtually unlimited and the interplay between formation criteria and long-term stability becomes paramount.